\documentclass[11pt,a4paper]{article}
\usepackage[margin=25mm]{geometry}
\usepackage[T1]{fontenc}
\usepackage{lmodern}
\usepackage{amsmath,amssymb}
\usepackage{booktabs,longtable,array}
\usepackage{graphicx}
\usepackage{hyperref}
\hypersetup{colorlinks=true,linkcolor=blue,urlcolor=blue}
\setlength{\parskip}{0.5em}
\setlength{\parindent}{0pt}
\setlength{\emergencystretch}{3em}
\newcommand{\starred}[1]{#1^{*}}
\newcommand{\NAIRU}{u^{*}}
\newcommand{\gap}{g}
\newcommand{\eps}{\varepsilon}

\title{Australian SMOG Model\\Coefficient and Fit Analysis}
\author{Model output review}
\date{16 September 2026\\Saved model results through 2026Q1}

\begin{document}
\maketitle

\section{Scope and interpretation}
This report analyzes the standalone Australian Python SMOG model.
Estimates come from its saved
parameter JSON and accepted-iteration likelihood history. The generation
sample is 1980Q1--2023Q4 (176 quarters). The update applies those fixed
coefficients to later data; the latest complete saved result is 2026Q1.
The latent potential output and NAIRU are model estimates, not directly
observed quantities.

\section{Results and charts}
The fixed-parameter update runs through 2026Q1. Its saved smoothed output
gap is $-0.371\%$ of potential output and its NAIRU is $4.089\%$.
Observed unemployment is $4.286\%$, or 0.197 percentage points above the
fitted NAIRU. Smoothed states use the full updated sample; filtered states
use observations available through each plotted quarter, with the same
historical coefficients. The charts are model estimates, not direct
measurements of potential output or the NAIRU.

\begin{figure}[p]
\centering
\includegraphics[width=0.85\linewidth,height=0.28\textheight,keepaspectratio]{\detokenize{../results/figures/au_smog_output_gap.png}}
\caption{Output gap: smoothed and filtered estimates.}
\end{figure}

\begin{figure}[p]
\centering
\includegraphics[width=0.85\linewidth,height=0.28\textheight,keepaspectratio]{\detokenize{../results/figures/au_smog_unemployment_and_nairu.png}}
\caption{Observed unemployment and smoothed and filtered NAIRU.}
\end{figure}

\begin{figure}[p]
\centering
\includegraphics[width=0.85\linewidth,height=0.28\textheight,keepaspectratio]{\detokenize{../results/figures/au_smog_unemployment_gap.png}}
\caption{NAIRU minus observed unemployment: smoothed and filtered estimates.}
\end{figure}

\clearpage

\section{Reading the current output gap}
The negative gap and unemployment above the NAIRU are consistent with the
negative unemployment-gap loading in this model. Yet the latest local inputs
also show trimmed inflation of $3.466\%$ and the model's constructed
expectation measure of $3.126\%$ in 2026Q1. In the fitted inflation equation,
the smoothed output gap contributes only $-0.014$ percentage points. The
other fitted terms produce a conditional inflation value of $3.294\%$ in
total; actual inflation is $0.173$ points higher. This is a decomposition
using the smoothed state, not an out-of-sample forecast.

It is reasonable to question whether the model's unemployment-driven gap
captures every source of current price pressure. The observation does not,
on its own, establish that the gap is too negative: lagged inflation,
expectations, import prices, and the inflation residual also enter the
measurement equation. The near-zero unemployment measurement-noise estimate
and the NAIRU innovation floor make the latent gap sensitive to the imposed
structure. Compare it with other capacity measures before treating its level
as a definitive estimate.

\section{State-space structure}
The model uses seven states,
$s_t=(\gap_t,\gap_{t-1},\gap_{t-2},y^*_t,y^*_{t-1},\NAIRU_t,\NAIRU_{t-1})'$, where
$y_t$ is log output, $\gap_t=y_t-y_t^*$ is a log output gap, and $\NAIRU_t$
is the latent structural unemployment rate. The reported output gap is
$100\gap_t$, in percent of potential output. The transition equations
are
\begin{align}
 \gap_t &= \phi_1\gap_{t-1}+\phi_2\gap_{t-2}+\eta^g_t,\\
 y_t^* &= y_{t-1}^*+g^*+\eta^y_t,\\
 \NAIRU_t &= \NAIRU_{t-1}+\eta^u_t.
\end{align}
The output gap follows an AR(2), potential log output is a random walk with
quarterly drift $g^*$, and the NAIRU is a random walk. The three state shocks
are independent in the implemented diagonal covariance matrix. The state
vector also carries lags needed by the observation equations.

At the beginning of each generation run, HP-based gap, potential output, and
NAIRU values set a known initial state with zero prior covariance. The model recomputes these HP series from its historical workbook. The HP series initialize the
filter and provide regression-based optimizer starting values; the saved
maximum-likelihood coefficients are fitted by Python from the country's own
signals. This fixed initialization can matter for early likelihood terms and
the resulting parameter estimates.

\section{Australia: equations and estimates}
Australia has three observed signals: log non-farm GDP $y_t$, unemployment
$u_t$, and trimmed inflation $\pi_t$. Let $\pi_t^e$ be the four-quarter moving
average of year-on-year inflation, $p_{t-1}^m$ the lagged four-quarter change
in the import price deflator, and $n_{t-1}$ lagged nominal unit labour cost
growth. Let $D_t=1$ after 1993Q1 and zero otherwise, and $C_t$ denote the
COVID dummy. The implemented observation equations are
\begin{align}
 y_t &= y_t^*+\gap_t+\phi_3 C_t+\eps^y_t,\\
 u_t &= \NAIRU_t+\lambda_3\gap_t
       +\rho(u_{t-1}-\NAIRU_{t-1})+\eps^u_t,\\
 \pi_t &= D_t\{(1-\beta_1-\beta_2)\pi_t^e
                  +\beta_1\pi_{t-1}+\beta_2\pi_{t-2}\}\notag\\
 &\quad +(1-D_t)\{(1-\beta_1-\beta_2-\beta_3-\gamma)\pi_t^e
              +\beta_1\pi_{t-1}+\beta_2\pi_{t-2}
              +\beta_3\pi_{t-3}+\gamma n_{t-1}\}\notag\\
 &\quad +\lambda_1\gap_t+\psi p_{t-1}^m+\eps^\pi_t.
\end{align}
The inflation expectation weight is therefore implied by the other
coefficients, separately in the two regimes. It is not an independent
parameter. Unit labour cost growth is a regressor in the pre-targeting
Australian Phillips curve, not an additional observed signal.

\begin{longtable}{@{}>{$}l<{$}r p{0.65\linewidth}@{}}
\caption{Australian maximum-likelihood coefficients. Values are rounded.}\\
\toprule
\text{Parameter} & \multicolumn{1}{c}{\text{Estimate}} & Meaning in the implemented model\\
\midrule
\endfirsthead
\toprule
\text{Parameter} & \multicolumn{1}{c}{\text{Estimate}} & Meaning\\
\midrule
\endhead
\beta_1 & -0.380964 & Weight on one-quarter lagged inflation in both regimes.\\
\beta_2 & -1.201930 & Weight on two-quarter lagged inflation in both regimes.\\
\beta_3 & -0.652679 & Weight on three-quarter lagged inflation before targeting.\\
\gamma & 0.004812 & Weight on lagged unit labour cost growth before targeting.\\
\lambda_1 & 3.739261 & Inflation response to a one-unit log output gap; a one percentage point gap gives about 0.0374 percentage points of inflation through this term.\\
\lambda_3 & -29.930362 & Unemployment response to the log output gap; a one percentage point gap corresponds to about $-0.2993$ percentage points of unemployment, holding other terms fixed.\\
\rho & 0.626289 & Persistence of lagged unemployment relative to the previous NAIRU in the unemployment equation.\\
\phi_1 & 1.295685 & First autoregressive coefficient of the latent output gap.\\
\phi_2 & -0.364960 & Second autoregressive coefficient of the latent output gap.\\
\phi_3 & -0.000494 & COVID dummy effect on observed log GDP, conditional on latent output.\\
\psi & 0.025443 & Loading of lagged four-quarter import-price change in inflation.\\
g^* & 0.007482 & Quarterly drift of potential log output, about 3.0\% at an annualized log-growth rate.\\
\eps_1 & 0.004307 & Standard deviation of log GDP measurement noise; about 0.431\% in log-output terms.\\
\eps_2 & 0.000010 & Standard deviation of unemployment measurement noise, in percentage points; at the imposed lower bound.\\
\eps_3 & 0.334839 & Standard deviation of inflation measurement noise, in percentage points.\\
\eps_5 & 0.005644 & Standard deviation of the log output-gap state shock; about 0.564\% of potential output.\\
\eps_6 & 0.004184 & Standard deviation of the potential-log-output state shock.\\
\eps_7 & 0.159770 & Standard deviation of the quarterly NAIRU state shock, in unemployment percentage points.\\
\bottomrule
\end{longtable}

The implied expectation weights are $2.5829$ after 1993Q1 and $3.2308$
before targeting. The large positive implied weights and negative lag weights
are features of this fitted equation, not general statements about inflation
dynamics. In particular, the near-zero fitted $\eps_2$ makes the unemployment
equation almost exact conditional on the latent states. The code constrains
$\eps_7$ to be at least half the historical standard deviation of quarterly
unemployment changes because the unrestricted likelihood allowed the NAIRU
innovation scale to collapse. This bound and the $\eps_2$ boundary solution
mean the NAIRU and output-gap decomposition should be treated as sensitive to
the model's restrictions.

\section{Likelihood fit and convergence}
The script uses SciPy's bounded L-BFGS-B optimizer on the Kalman-filter log
likelihood. Starting guesses come from Australian historical data. The
accepted-iteration history stores all 18 parameters, the full-sample log
likelihood, and quarter-level contributions.

\begin{center}
\begin{tabular}{@{}lr@{}}
\toprule
Fit statistic & Australia\\
\midrule
Historical quarters & 176\\
Accepted iterations & 109\\
Initial log likelihood & 234.1271\\
Iteration 25 & 549.1910\\
Iteration 50 & 549.2789\\
Final log likelihood & 549.2792\\
\bottomrule
\end{tabular}
\end{center}

Accepted steps increased the likelihood monotonically. The last history row
agrees with the saved parameter JSON. Numerical convergence from these
starting values and bounds does not establish a global optimum, tight
identification, or forecast accuracy. The script does not export standard
errors or out-of-sample forecast scores. The quarter-level likelihood file
can identify observations that were hardest to predict: each contribution is
a Kalman one-step-ahead predictive density term, not a squared-error metric.
Early missing regressors cause the affected signals to be masked.

\section{Reproduction and source files}
From this model folder, run the generation script to re-estimate historical
coefficients and overwrite the history CSVs. Then run the update script to
apply the fixed fit to later quarters. Source files used for this report:
\begin{itemize}
\item Generation and update scripts in the \texttt{code} folder.
\item Fitted parameter JSON at the model-folder root.
\item The fit-history and likelihood-by-quarter CSVs in \texttt{output}.
\item The updated state CSV in \texttt{results} and the \texttt{Update}
      tab of the local model-input workbook.
\end{itemize}
Numeric tables are a snapshot and do not update automatically when Python is
rerun.

\end{document}
